\subsection{2.13 Projektionsprinzip als zentrales Paradigma der PST}

Mit den Ergebnissen der numerischen Untersuchungen hat sich die Perspektive der
Projective Structure Theory (PST) grundlegend verändert.

Zu Beginn der Entwicklung stand die Frage:

\begin{quote}
\emph{„Besitzt der Ontophänomenale Phasenraum (OPP) selbst bereits eine
vierdimensionale Struktur?“}
\end{quote}

Nach zahlreichen Simulationen lautet die heutige Antwort:

\[
\boxed{
\text{Wahrscheinlich nicht.}
}
\]

Stattdessen verdichten sich die Hinweise darauf, dass die beobachtete Physik
erst durch einen geeigneten Projektionsprozess entsteht.

Damit verschiebt sich der Schwerpunkt der Theorie vollständig.

%%%%%%%%%%%%%%%%%%%%%%%%%%%%%%%%%%%%%%%%%%%%%%%%%%%%%%%%%%%%%%%

\subsubsection*{2.13.1 Das klassische Paradigma}

In den meisten physikalischen Theorien wird implizit angenommen:

\[
\boxed{
\text{Fundamentale Realität}
=
\text{Raumzeit}
}
\]

Alle weiteren Strukturen

\begin{itemize}
\item Materie,
\item Felder,
\item Wechselwirkungen,
\item Quantenzustände
\end{itemize}

leben innerhalb dieser Raumzeit.

Raum und Zeit werden somit als ontologisch fundamental betrachtet.

%%%%%%%%%%%%%%%%%%%%%%%%%%%%%%%%%%%%%%%%%%%%%%%%%%%%%%%%%%%%%%%

\subsubsection*{2.13.2 Das PST-Paradigma}

Die PST schlägt genau die umgekehrte Sichtweise vor.

Nicht Raumzeit ist fundamental.

Fundamental ist ausschließlich der OPP.

Die Kette lautet daher

\[
\boxed{
\text{OPP}
\longrightarrow
\text{Projektion}
\longrightarrow
\text{Raumzeit}
}
\]

Raumzeit besitzt damit denselben ontologischen Status wie beispielsweise

\begin{itemize}
\item Temperatur,
\item Druck,
\item Schall,
\item Flüssigkeit.
\end{itemize}

Alles dies sind emergente Beschreibungen tieferliegender Strukturen.

%%%%%%%%%%%%%%%%%%%%%%%%%%%%%%%%%%%%%%%%%%%%%%%%%%%%%%%%%%%%%%%

\subsubsection*{2.13.3 Projektion statt Einbettung}

Ein wichtiger Unterschied zu vielen höherdimensionalen Theorien besteht darin,
dass die PST keine Einbettung einer 4D-Raumzeit in einen höheren Raum fordert.

Sie fordert stattdessen einen Projektionsoperator

\[
\mathcal P.
\]

Formal

\[
\mathcal P :
\mathrm{OPP}
\rightarrow
\mathcal M.
\]

Dabei bezeichnet

\[
\mathcal M
\]

die beobachtbare physikalische Welt.

Diese Projektion reduziert Information.

Nicht alle Freiheitsgrade des OPP erscheinen innerhalb der Physik.

%%%%%%%%%%%%%%%%%%%%%%%%%%%%%%%%%%%%%%%%%%%%%%%%%%%%%%%%%%%%%%%

\subsubsection*{2.13.4 Informationsverlust}

Jede Projektion besitzt die Eigenschaft

\[
\dim(\mathcal M)
<
\dim(\mathrm{OPP}).
\]

Dabei geht Information verloren.

Formal

\[
I(\mathcal M)
<
I(\mathrm{OPP}).
\]

Dieser Informationsverlust ist kein Fehler.

Er ist gerade der Mechanismus,

durch den stabile makroskopische Strukturen entstehen.

%%%%%%%%%%%%%%%%%%%%%%%%%%%%%%%%%%%%%%%%%%%%%%%%%%%%%%%%%%%%%%%

\subsubsection*{2.13.5 Warum Projektionen überhaupt?}

Die numerischen Untersuchungen liefern hierfür erstmals eine konkrete Motivation.

Der rohe OPP besitzt

\begin{itemize}
\item keine stabile Vierdimensionalität,
\item keine Lorentzstruktur,
\item keine klar erkennbare Raumzeit.
\end{itemize}

Erst geeignete Projektionen erzeugen

\[
d_{\mathrm{eff}}
\approx
4.
\]

Damit entsteht erstmals eine plausible Erklärung,

warum unsere Welt gerade vierdimensionale Eigenschaften besitzt.

%%%%%%%%%%%%%%%%%%%%%%%%%%%%%%%%%%%%%%%%%%%%%%%%%%%%%%%%%%%%%%%

\subsubsection*{2.13.6 Beobachter als Projektionsoperator}

Eine besonders interessante Konsequenz ergibt sich,

wenn Beobachter selbst als Projektionsprozesse verstanden werden.

Dann gilt

\[
\boxed{
\text{Beobachtung}
=
\text{Projektion}.
}
\]

Ein physikalischer Beobachter besitzt niemals Zugang zum vollständigen OPP.

Er registriert lediglich

\[
\mathcal P(\mathrm{OPP}).
\]

Damit wird verständlich,

warum alle Messungen grundsätzlich informationsbegrenzt sind.

%%%%%%%%%%%%%%%%%%%%%%%%%%%%%%%%%%%%%%%%%%%%%%%%%%%%%%%%%%%%%%%

\subsubsection*{2.13.7 Verbindung zur Quantisierung}

Diese Sichtweise eröffnet gleichzeitig einen neuen Zugang zur Quantisierung.

Bisher wird häufig angenommen,

dass Quantisierung eine fundamentale Eigenschaft der Natur sei.

Die PST schlägt stattdessen vor,

dass Quantisierung eine Folge endlicher Projektionsauflösung ist.

Anschaulich:

Ein kontinuierlicher OPP

\[
\longrightarrow
\]

wird durch endliche Auflösung beobachtet

\[
\longrightarrow
\]

es entstehen diskrete Zustände.

Somit wäre Quantisierung kein ontologisches Fundament,

sondern eine Konsequenz des Projektionsprozesses.

%%%%%%%%%%%%%%%%%%%%%%%%%%%%%%%%%%%%%%%%%%%%%%%%%%%%%%%%%%%%%%%

\subsubsection*{2.13.8 Verbindung zur Raumzeit}

Genau dieselbe Argumentation lässt sich nun auf Raum und Zeit anwenden.

Nicht nur Quantisierung entsteht projektionsbedingt.

Auch Raumzeit selbst könnte eine Folge endlicher Beobachtbarkeit sein.

Formal:

\[
\boxed{
\text{Raum}
=
\mathcal P_R(\mathrm{OPP}),
}
\]

\[
\boxed{
\text{Zeit}
=
\mathcal P_T(\mathrm{OPP}).
}
\]

Damit erhalten Raum und Zeit denselben Status wie Quantisierung.

Beides sind Eigenschaften der Projektion,

nicht Eigenschaften des OPP.

%%%%%%%%%%%%%%%%%%%%%%%%%%%%%%%%%%%%%%%%%%%%%%%%%%%%%%%%%%%%%%%

\subsubsection*{2.13.9 Verbindung zur Relativität der Existenz}

Im ersten Band wurde die Relativität der Existenz (RdE) eingeführt.

Sie besagt,

dass Bestimmtheit beobachterabhängig ist.

Dies fügt sich nun nahtlos ein.

Denn unterschiedliche Beobachter besitzen unterschiedliche Projektionen.

Formal

\[
\mathcal P_A
\neq
\mathcal P_B.
\]

Daher gilt allgemein

\[
\boxed{
\mathcal P_A(\mathrm{OPP})
\neq
\mathcal P_B(\mathrm{OPP}).
}
\]

Die Relativität der Existenz erscheint somit als unmittelbare Konsequenz
unterschiedlicher Projektionsoperatoren.

%%%%%%%%%%%%%%%%%%%%%%%%%%%%%%%%%%%%%%%%%%%%%%%%%%%%%%%%%%%%%%%

\subsubsection*{2.13.10 Offene mathematische Fragen}

Die numerischen Arbeiten beantworten jedoch nicht,

wie der physikalisch korrekte Projektionsoperator aussieht.

Insbesondere bleiben offen:

\begin{enumerate}

\item
Warum erzeugt Diffusion Maps so stabile Ergebnisse?

\item
Warum funktioniert Kernel PCA fast ebenso gut?

\item
Existiert ein universeller Operator

\[
\mathcal P^\ast,
\]

aus dem beide Verfahren als Approximation hervorgehen?

\item
Kann dieser Operator aus den Axiomen des OP hergeleitet werden?

\item
Erzeugt genau dieser Operator automatisch

\[
SO(1,3),
\quad
SU(3),
\quad
SU(2),
\quad
U(1)?
\]

\end{enumerate}

Diese Fragen bilden gegenwärtig den mathematischen Kern der PST.

%%%%%%%%%%%%%%%%%%%%%%%%%%%%%%%%%%%%%%%%%%%%%%%%%%%%%%%%%%%%%%%

\subsubsection*{2.13.11 Das neue Forschungsprogramm}

Die bisherige Forschung konzentrierte sich hauptsächlich auf den OPP.

Die zukünftige Forschung verschiebt den Schwerpunkt.

Gesucht wird nicht länger primär die Struktur des OPP,

sondern die Struktur der Projektion.

Das zukünftige Forschungsprogramm lautet daher:

\[
\boxed{
\text{Axiome}
\rightarrow
\text{OPP}
\rightarrow
\boxed{\mathcal P}
\rightarrow
\text{Raumzeit}
\rightarrow
\text{Symmetrien}
\rightarrow
\text{Standardmodell}.
}
\]

Der Projektionsoperator wird damit zum eigentlichen Herzstück der PST.

%%%%%%%%%%%%%%%%%%%%%%%%%%%%%%%%%%%%%%%%%%%%%%%%%%%%%%%%%%%%%%%

\subsubsection*{2.13.12 Schlussfolgerung}

Die numerischen Untersuchungen markieren einen entscheidenden Wendepunkt.

Sie liefern erstmals konkrete Hinweise darauf,

dass die vierdimensionale Raumzeit möglicherweise gar nicht die fundamentale
Struktur der Wirklichkeit ist.

Fundamental wäre stattdessen ein wesentlich reichhaltigerer OPP,

dessen physikalisch beobachtbare Eigenschaften erst durch geeignete
Projektionsprozesse entstehen.

Damit entwickelt sich die PST von einer Theorie des OPP zunehmend zu einer
Theorie der Projektion.

Dieser Perspektivwechsel dürfte sich als einer der zentralen theoretischen
Fortschritte der bisherigen Entwicklung erweisen.

